%% main.tex — Measured EZKL circuit size of transformers: what costs rows, and what it costs in perplexity

\documentclass[10pt]{article}

%% --- Page geometry ---
\usepackage[margin=1in]{geometry}

%% --- Core packages ---
\usepackage[utf8]{inputenc}
\usepackage[T1]{fontenc}
\usepackage{amsmath, amssymb, amsthm}
\usepackage{mathtools}
\usepackage{bm}
\usepackage{booktabs}
\usepackage{graphicx}
\usepackage{hyperref}
\usepackage{xcolor}
\usepackage{natbib}
\usepackage{url}

%% --- Math macros ---
\newcommand{\dmodel}{d_{\mathrm{model}}}
\newcommand{\nlayers}{n_{\mathrm{layers}}}
\newcommand{\ZKcost}{\mathrm{ZK\_cost}}

%% ============================================================
\begin{document}

\title{What Costs Rows and Proving Time in Small EZKL Transformer Circuits,\\
       and What It Costs in Perplexity}

\author{%
  Nissan Pow\\
  Orchestral Labs
}

\maketitle

% ============================================================
\begin{abstract}
% ============================================================
Design advice for zero-knowledge (ZK) transformer inference rests on analytical
constraint-count models and on swapping activations for cheaper ones. We test an
analytical model against EZKL~23.0.5 on small causal transformers and measure
what each choice costs in perplexity. We compile 108 configurations (width
64--128, 1--4 layers, sequence length 8--32, four activation variants) and
generate 24 verified proofs. The analytical model ranks compiled row counts well
(Spearman $\rho{=}0.93$) but no better than the product of the dimensions ($0.96$),
is uncalibrated (measured/modelled rows $0.12$--$0.53$), and predicts held-out
widths, depths and sequence lengths worse ($24$--$38\%$ error) than a power law
($6$--$18\%$) or a calibrated additive model ($8$--$11\%$). On 16 out-of-range configurations (width 256--512) the analytical model's
fitted prediction errs by $122\%$ against $10\%$ for the dimension product, and
the activation overheads persist (GELU $1.45$--$1.54\times$, ReLU
$1.08$--$1.17\times$ the linear variant's rows). Rows predict proving
time (Spearman $0.89$ over six shapes, no better than the dimension product), and
proving time is a step function of the power-of-two proving domain: crossing a
boundary doubled it. Repeated runs on a shared machine differ by up to $1.9\times$,
so proving times do not separate the nonlinearities. For
sequences up to 32, the FFN nonlinearity dominates the row overhead for GELU and poly2 (at width
128, ReLU $+7\%$, poly2 $+20\%$, GELU $+42\%$), but softmax grows with sequence length
(from $4\%$ to $24\%$ at width 64 as the length goes from 8 to 64). EZKL's domain
size also depends on an implementation detail, the causal-mask fill value, which
moved it from 19 to 26 at constant rows. In a WikiText-2 study (3 seeds,
validation-selected checkpoints), ReLU, poly2 and GELU are within $1\%$ in
perplexity, larger models help, and nothing trained from scratch approaches
pretrained GPT-2. Compiled circuits reach width 512 (the grid width 128) with at most four layers,
proofs only domains up to $2^{21}$, proving times come from one or two runs on a
shared machine, and quantized
accuracy is not measured. Code and
raw records are released.
\end{abstract}

% ============================================================
\section{Introduction}
\label{sec:intro}
% ============================================================

Verifiable inference---proving that a model ran correctly on private
inputs---needs a ZK proof of the network's forward pass, and proving
transformer-scale models is costly~\citep{zkllm2024,zkgpt2025,nanozk2026}. A
common design heuristic substitutes polynomial activations for GELU and softmax;
an analytical counter-argument says matrix multiplication dominates the circuit
so activation choice barely matters and shape ($\dmodel$, $\nlayers$) is the
lever. Both are statements about cost models, not measurements. We measure, for
one prover version, (i) how well an analytical constraint model predicts compiled
circuit size, (ii) which architectural choices actually cost rows, and (iii) what
those choices cost in language-model quality. Our contributions:

\begin{enumerate}
  \item \textbf{A measured test of an analytical cost model}
    (Section~\ref{sec:ezkl}): 108 EZKL compilations, with trivial baselines, a
    calibrated additive competitor, a cluster bootstrap over shapes, and
    leave-one-axis-out tests; and 24 verified proofs relating rows to proving time,
    memory and key size (Section~\ref{sec:proving}).
  \item \textbf{An attention/FFN decomposition} of the activation overhead across
    sequence lengths (Section~\ref{sec:ezkl}), and the finding that EZKL's
    proving-domain size depends on the causal-mask fill value
    (Section~\ref{sec:domain}).
  \item \textbf{A leak-free, checkpoint-selected quality study}
    (Section~\ref{sec:quality}) and a joint reading of cost and quality
    (Section~\ref{sec:joint}).
\end{enumerate}

\paragraph{What this paper does not show.}
We prove only 24 small circuits (domain at most $2^{21}$) from one or two noisy runs each; we do
not measure circuits larger than width 128; and we do not study other proof
systems (Section~\ref{sec:limits}).

% ============================================================
\section{Related Work}
\label{sec:related}
% ============================================================

\paragraph{ZK-ML proof systems.}
\citet{zkllm2024} introduce tlookup and zkAttn arguments for transformer
nonlinearities; \citet{zkgpt2025} reduce the cost of non-linear layers with
constraint fusion; \citet{nanozk2026} prove each layer separately and chain the
layer proofs with commitments; \citet{telescarse2025} shrink circuits by
sparsification and by narrowing activation ranges, which shrinks the lookup
tables; \citet{kang2022scaling} proved inference of an ImageNet-scale model with
ZK-SNARKs. These works fix the architecture and improve the prover or the
model. We hold the prover (EZKL~\citep{ezkl}) fixed and measure how architecture
choices change compiled circuit size and proving cost.

\paragraph{Hardware-aware NAS and cost models.}
\citet{cai2019proxylessnas}, \citet{tan2019mnasnet} and \citet{wu2019fbnet}
established latency-aware search with hardware cost proxies, validated against
measured latency; \citet{javaheripi2022lts} observe that decoder parameter count
correlates with perplexity and use it as a training-free proxy for language
models. \citet{nass2020} search for architectures under homomorphic-encryption and
garbled-circuit secure inference. Our analogue of the latency-validation step is
Sections~\ref{sec:ezkl} and~\ref{sec:proving}, with circuit size and proving time
standing in for latency.

\paragraph{Zero-cost proxies.}
SNIP~\citep{lee2019snip} and NEAR~\citep{husistein2025near} are training-free
proxies for network quality; \citet{serianni2023understanding} study
training-free metrics for RNNs and transformers. We do not evaluate proxies.

% ============================================================
\section{Setup}
\label{sec:setup}
% ============================================================

\paragraph{Models.}
Decoder-only transformers with pre-LayerNorm, tied embeddings, causal attention,
and an FFN activation from \{GELU, ReLU, poly2 ($x^2$), identity\}. Attention is
either softmax or unnormalized (``linear'') causal attention. The variant called
``linear'' throughout removes \emph{both} the softmax and the FFN nonlinearity
(identity FFN); Sections~\ref{sec:ezkl} and~\ref{sec:quality} also report the two
factors separately. Positions are learned absolute embeddings unless rotary is
stated; the EZKL grid uses learned absolute positions.

\paragraph{Analytical cost model.}
The model (\path{experiments/operator_attribution.py}; per-operator costs vendored
in \path{proxy/op_costs.py}) sums, per operator category, an operation count
times a per-operation constraint cost. Matrix products cost 1 constraint per
multiply-accumulate, counted as 2 MACs per product: per layer, QKV and output
projections $8sd^2$, attention scores and values $4s^2d$, and FFN up and down
projections $16\,sd^2$ at an expansion of 4. Softmax costs 32 per score
element, the FFN activation costs 2 (ReLU), 24 (GELU), 2 (poly2) or 0
(identity) per hidden element, LayerNorm 24 per element, and the token embedding
(plus position embedding when rotary is off) 2 per table entry. It has no term
for the output (LM-head) projection. Here $s$ is sequence length and $d$ the
width. We call its output ``modelled constraints.'' It is our own formalization of the
constraint-counting argument (the code is in the repository), a stand-in for the
heuristics common in the ZK-ML literature, not a published system's model.

\paragraph{Measurement.}
Each model is exported to ONNX (opset 18) and passed to EZKL~23.0.5
(\texttt{gen\_settings}, then \texttt{calibrate\_settings} with
\texttt{target=resources}). Calibration uses four fixed random token sequences
and chooses the fixed-point scale itself (\texttt{param\_scale} is 13 in 86
records, 12 in 17 and 11 in 5; the linear variant is always 13), which
confounds row counts with the scale, so we replicate the grid with the scale
pinned (Section~\ref{sec:ezkl}, ``Robustness to the fixed-point scale''). We record \texttt{num\_rows} and \texttt{logrows}.
The grid compiles only (circuit size, not proving time); Section~\ref{sec:proving}
proves a subset. Weights are random initialization. The causal mask uses a finite fill value of
$-100$ (Section~\ref{sec:domain}). \texttt{total\_assignments} equals $2\times$
\texttt{num\_rows} (exactly in 72 of 108 records, within $10^{-4}$ in the rest), so
we report rows only.

% ============================================================
\section{Measured EZKL Circuit Size}
\label{sec:ezkl}
% ============================================================

The grid is the full product of width $\in\{64,96,128\}$, layers
$\in\{1,2,4\}$, sequence length $\in\{8,16,32\}$ and four activation variants
(108 configurations; vocabulary 512; 4 heads). All 108 compiled. The grid has
only 27 distinct shapes, and the four variants at a shape are near-duplicates, so
all intervals below are cluster bootstraps over shapes.

\paragraph{The analytical model ranks well, but not better than a trivial baseline.}
The Spearman correlation between modelled constraints and measured rows is
$0.927$ (95\% CI $[0.83,0.97]$). The product $\dmodel\nlayers s$ achieves $0.964$
($[0.91,0.98]$) and $\nlayers s$ alone $0.935$; the paired difference
(analytical minus product) has 95\% CI $[-0.091,+0.005]$. Within each variant the
analytical model scores $0.92$--$0.94$ and the product $0.97$--$0.98$. On this
grid, a monotone size proxy ranks configurations at least as well as the model.

\paragraph{The model is uncalibrated, and it extrapolates worse than every fitted alternative.}
The ratio of measured rows to modelled constraints ranges from $0.12$ to $0.53$
(median $0.22$), so rows are \emph{smaller} than the modelled count and no single
constant converts one to the other. A log-log fit on the model output alone has
slope $0.77$ and $R^2{=}0.85$; a fit on the three dimensions gives
$\mathrm{rows}\propto \dmodel^{0.71}\nlayers^{0.79}s^{1.10}$ with $R^2{=}0.97$.
The model also misorders $27$ of $162$ within-shape variant pairs.

To test transfer, Table~\ref{tab:loo} holds out one value of one axis (each of
three widths, three depths or three sequence lengths in turn; the smallest and
largest values are extrapolations, the middle one an interpolation), fits on the
remaining configurations, and predicts the held-out rows (mean absolute
percentage error, averaged over held-out values). Besides the model and the
product $\dmodel\nlayers s$, it includes a three-exponent power law and a
\emph{calibrated additive model}: non-negative least squares of rows on the
model's own operator categories (matrix products, softmax, activation,
LayerNorm, embedding), optionally with an output-projection term the model
lacks. The analytical model with a fitted log-log line is worst on every axis; the
additive fit is best on width, within $1$ point of the best on depth and $2$
points behind the power law on sequence length.

\begin{table}[t]
\centering
\caption{Leave-one-axis-out error of measured rows (mean absolute percentage
  error, 108 configurations; lower is better). Log-log fits include
  per-variant offsets; the additive fits have 5 or 6 non-negative weights.}
\label{tab:loo}
\begin{tabular}{lccc}
\toprule
Predictor & held-out $\dmodel$ & held-out $\nlayers$ & held-out $s$ \\
\midrule
Analytical model (log-log fit) & 38.1 & 24.0 & 37.0 \\
Product $\dmodel\nlayers s$ (log-log fit) & 14.7 & 17.2 & 21.4 \\
Power law $\dmodel^{x}\nlayers^{y}s^{z}$ & 17.6 & 7.2 & 6.0 \\
Additive NNLS, model's categories & 11.2 & 7.5 & 8.0 \\
Additive NNLS $+$ LM-head term & 11.8 & 7.6 & 8.1 \\
\bottomrule
\end{tabular}
\end{table}

The fitted additive weights (rows per modelled constraint) are $0$ for the matrix
product and embedding terms, $2.7$ for softmax, $1.4$ for the activation and
$8.0$ for LayerNorm; they exceed one because the matrix-product term, which
dominates the modelled total, receives no weight. Under a cluster bootstrap over
shapes ($1000$ resamples) the matrix-product weight is exactly zero in $92\%$ of
fits (95\% interval $[0,0.011]$). This does not show that matrix products are
free: the matrix-product and LayerNorm terms are strongly collinear (Spearman
$0.97$), the fit without the matrix-product term explains the same variance as
the full fit ($R^2{=}0.989$ in both), and the matrix-product term alone gives
$R^2{=}0.79$. The data cannot separate the two terms, so we read the result only
as: a model that charges for matrix products in proportion to multiply-accumulate
counts does not describe EZKL's row counts, consistent with EZKL packing many
multiply-accumulates into one row.

\paragraph{What costs rows: the FFN nonlinearity at short sequences.}
Table~\ref{tab:decomp} crosses attention (softmax or linear) with the FFN
activation at $\nlayers{=}2$, $s{=}16$, and reports rows relative to linear
attention with an identity FFN.

\begin{table}[t]
\centering
\caption{Measured EZKL rows relative to linear attention with identity FFN,
  same shape ($\nlayers{=}2$, $s{=}16$). Entries are the value at width 128 with
  the range over widths 64--128 in parentheses.}
\label{tab:decomp}
\begin{tabular}{lcc}
\toprule
FFN activation & Linear attention & Softmax attention \\
\midrule
identity & 1.00 & 1.02 (1.02--1.08) \\
ReLU     & 1.07 (1.06--1.07) & 1.09 (1.09--1.14) \\
poly2    & 1.20 (1.17--1.22) & 1.23 (1.23--1.25) \\
GELU     & 1.42 (1.36--1.44) & 1.44 (1.44--1.48) \\
\bottomrule
\end{tabular}
\end{table}

Softmax adds $2$--$8\%$ (falling with width: $1.079$, $1.031$, $1.022$ at widths
$64$, $96$, $128$); the FFN activation adds $7\%$ (ReLU), $20\%$ (poly2) or
$42\%$ (GELU). Over the full grid, relative to the identity-FFN linear variant,
GELU uses $1.45\times$ the rows (by width $1.43$, $1.46$, $1.43$), poly2
$1.25\times$ ($1.26$, $1.24$, $1.22$) and ReLU $1.10\times$ ($1.14$, $1.10$,
$1.09$); GELU's overhead is flat in width, while poly2's and ReLU's shrink slightly
(by 3 and 5 points between widths 64 and 128). The analytical model puts the linear variant only $3$--$8.5\%$ below
GELU at these shapes, against $31\%$ measured.

\paragraph{Softmax grows with sequence length.}
Table~\ref{tab:seqT} repeats the decomposition at sequence lengths $8$--$64$
for widths 64 and 128 (two layers, scale pinned at 11). The softmax overhead
grows quickly with sequence length and falls with width: at width 64 it rises
from $4\%$ ($s{=}8$) to $24\%$ ($s{=}64$), where it equals the GELU overhead
($23\%$ over softmax attention with identity FFN); at width 128 it rises from
$1.3\%$ to $6.3\%$. The FFN nonlinearity therefore dominates only at short
sequences and wide models; at width 64 and $s{=}64$, softmax is as costly as
GELU, and because softmax is quadratic in $s$ it should dominate at the longer
sequences used in practice, which we did not measure. Across the whole
decomposition, ReLU uses $16$--$25\%$ fewer rows than GELU (and $16$--$29\%$
across the 27 grid shapes, median $23\%$).

\begin{table}[t]
\centering
\caption{Measured rows relative to linear attention with identity FFN at the same
  shape, versus sequence length ($\nlayers{=}2$, scale pinned at 11). All
  entries use softmax attention.}
\label{tab:seqT}
\begin{tabular}{llcccc}
\toprule
Width & FFN activation & $s{=}8$ & $s{=}16$ & $s{=}32$ & $s{=}64$ \\
\midrule
64  & identity & 1.04 & 1.08 & 1.14 & 1.24 \\
64  & ReLU     & 1.11 & 1.14 & 1.20 & 1.29 \\
64  & GELU     & 1.42 & 1.44 & 1.47 & 1.52 \\
128 & identity & 1.01 & 1.02 & 1.04 & 1.06 \\
128 & ReLU     & 1.09 & 1.09 & 1.10 & 1.12 \\
128 & GELU     & 1.45 & 1.44 & 1.43 & 1.41 \\
\bottomrule
\end{tabular}
\end{table}

\paragraph{Robustness to the fixed-point scale.}
Calibration chose the scale in the main grid (11--13). Re-running all 108
configurations with the scale pinned at 11 gives identical \texttt{logrows} in
every configuration and row counts that are identical in 72 of 108 and differ by
a few rows (below $3\times10^{-5}$ relative) in the rest, so every statistic above is
unchanged at the reported precision. (Weights are random and unseeded, which
accounts for the few-row differences between repeated compilations.)

\paragraph{Out-of-range check.}
To test extrapolation beyond the fitted range, we compiled 16 further
configurations (widths 256 and 512, 1--2 layers, sequence lengths 16 and 64;
scale pinned at 11; up to $2.2\times10^{7}$ rows at domain $2^{25}$) and predicted
their rows with every predictor fitted on the 108-configuration grid
(Table~\ref{tab:oor}). The analytical model's fitted line overpredicts by a factor of
about $2$ (error $122\%$), and its raw measured/modelled ratio falls to
$0.03$--$0.10$ (median $0.057$), below the $0.12$--$0.53$ seen on the grid: the gap
widens with size. The product $\dmodel\nlayers s$ and the power law predict
within about $10\%$; the additive fit, the best predictor in the leave-one-out
test, is worse out of range ($16\%$). Activation overheads persist at the larger
widths: relative to the linear variant, GELU uses $1.45$--$1.54\times$ the rows,
poly2 $1.23$--$1.31\times$ and ReLU $1.08$--$1.17\times$ (against $1.43$--$1.46$,
$1.22$--$1.26$ and $1.09$--$1.14$ at widths 64--128).

\begin{table}[t]
\centering
\caption{Prediction of rows for 16 out-of-range configurations (width 256--512,
  sequence length 16--64) by predictors fitted on the 108-configuration grid
  (width 64--128).}
\label{tab:oor}
\begin{tabular}{lccc}
\toprule
Predictor & Error (\%) & Median pred./measured & Spearman \\
\midrule
Analytical model (log-log fit) & 122.2 & 2.13 & 0.88 \\
Product $\dmodel\nlayers s$ (log-log fit) & 10.0 & 1.00 & 0.99 \\
Power law $\dmodel^{x}\nlayers^{y}s^{z}$ & 10.7 & 0.90 & 0.99 \\
Additive NNLS & 16.1 & 1.15 & 0.96 \\
\bottomrule
\end{tabular}
\end{table}

\paragraph{Scope.}
The grid has at most four layers, width 128 and sequence length 32, and the
out-of-range check reaches width 512 and sequence length 64 with 1--2 layers. The
analytical argument concerns much larger models (GPT-2 scale: width 768, 12
layers, sequence length 512 or more), which we do not reach; the power-law fits
should not be extrapolated beyond the range tested.

% ============================================================
\section{Proving Time, Memory and Key Size}
\label{sec:proving}
% ============================================================

Compiled rows are a proxy; to test it we generated real proofs for 24
configurations: six shapes (width 64--128, 1--2 layers, sequence length 8--16)
crossed with the four variants, at fixed-point scale 11 and domains $2^{19}$ to
$2^{21}$. Each configuration runs the full pipeline (compile, setup, witness,
prove, verify) in a fresh process, with a locally generated structured reference
string (suitable for timing, not for security). We ran the 20 configurations with
domain at most $2^{20}$ twice (the four at $2^{21}$ once). All proofs verified.

\begin{table}[t]
\centering
\caption{Mean proving time in seconds over two runs (domain $\log_2$ size in
  parentheses) for the four variants at six shapes (width/layers/sequence); the
  64/2/16 row is a single run. Shared 24-core machine; see the text for
  run-to-run variation.}
\label{tab:prove}
\begin{tabular}{lcccc}
\toprule
Shape & linear & ReLU & poly2 & GELU \\
\midrule
64/1/8   & 65 (19)  & 94 (19)  & 84 (19)  & 77 (19) \\
64/1/16  & 128 (20) & 166 (20) & 161 (20) & 156 (20) \\
64/2/8   & 141 (20) & 150 (20) & 141 (20) & 155 (20) \\
64/2/16  & 218 (21) & 253 (21) & 269 (21) & 346 (21) \\
96/1/8   & 67 (19)  & 96 (19)  & 76 (19)  & 146 (20) \\
128/1/8  & 138 (20) & 177 (20) & 160 (20) & 170 (20) \\
\bottomrule
\end{tabular}
\end{table}

\paragraph{Rows predict proving cost, through the domain size.}
Over the 24 configurations (six shapes, three domain levels), compiled rows
correlate with mean proving time (Spearman $0.89$, cluster-bootstrap 95\%
interval over the six shapes $[0.54,0.98]$); the product $\dmodel\nlayers s$
scores $0.88$. Within a shape, a larger row count means a longer mean proving time
in $26$ of $36$ variant pairs. Proving time is close to a step function of the
power-of-two domain: it roughly doubles with each extra domain bit ($65$--$96$\,s at
$2^{19}$, $128$--$177$\,s at $2^{20}$, and $218$--$346$\,s, single runs, at $2^{21}$).
Small row differences can therefore have a large effect: at width 96, GELU's
$552$K rows cross $2^{19}=524$K, its domain grows from 19 to 20 bits, and it takes
$146$\,s against $67$--$96$\,s for the other variants ($2.2\times$ the linear
variant). Even these toy models are heavy: the proving key is $8.5$--$40.5$\,GB,
peak memory $21$--$103$\,GB (decimal), the proof $0.7$--$1.3$\,MB, verification
$0.6$--$3.4$\,s, and (first run) setup $58$--$321$\,s.

\paragraph{Timing is noisy on a shared machine.}
The machine was shared with other jobs (load average $20$--$40$ on 24 cores). The
second run of each configuration took between $0.82\times$ and $1.91\times$ the
first run's time (median $1.24\times$); the per-configuration coefficient of
variation over the two runs is $15\%$ (median, maximum $44\%$). Absolute times,
and differences between variants of less than about $30\%$, should be read as
indicative only; the domain steps and the row-driven ordering across shapes are
the robust effects.

\paragraph{Activation overhead in proving time.}
Using mean times, relative to the linear variant ReLU takes a median $1.29\times$
as long to prove ($1.06$--$1.46$), poly2 $1.20\times$ ($1.00$--$1.30$) and GELU
$1.22\times$ ($1.10$--$2.19$, the maximum being the domain-boundary case), against
median row ratios of $1.09$, $1.20$ and $1.38$. Nonlinear activations therefore
cost roughly $20$--$30\%$ more proving time than the linear variant, but proving
time does not separate the nonlinearities: ReLU is faster than poly2 in one of six
shapes and slower in five, and GELU is slower than ReLU in three and slower than
poly2 in four. With two runs per configuration, these orderings are within the
run-to-run noise.

% ============================================================
\section{Proving-Domain Size Depends on the Mask Fill Value}
\label{sec:domain}
% ============================================================

EZKL's \texttt{logrows} sets the proving domain and hence key size and memory.
A causal softmax needs a masked-score fill value; $-\infty$ cannot be quantized
(\texttt{gen\_settings} fails), so a large negative constant is used.
Table~\ref{tab:fill} shows that this choice alone moves \texttt{logrows} while
the row count is unchanged: EZKL sizes its softmax lookup range from the most
negative input, so a fill of $-10^4$ gives a table spanning about $1.6\times10^8$
entries.

\begin{table}[t]
\centering
\caption{\texttt{logrows} versus causal-mask fill value at two shapes. Row counts
  are constant across fills (to within $2\times10^{-5}$, from unseeded random
  weights): $458{,}765$ (GELU), $375{,}821$ (ReLU), $407{,}565$ (poly2) and
  $347{,}629$ (linear) at width 64, 1 layer, $s{=}8$; $2.72$M, $2.05$M, $2.31$M
  and $1.88$M at width 128, 2 layers, $s{=}16$. The linear variant has no softmax, so no
  fill value applies and its single measurement is repeated. The
  calibration-chosen scale varied between 11 and 13 across cells; the domain sizes
  are set by the lookup range and are unaffected.}
\label{tab:fill}
\begin{tabular}{llccccc}
\toprule
Shape & Variant & $-10^4$ & $-10^3$ & $-100$ & $-30$ & $-10$ \\
\midrule
64/1/8 & GELU, ReLU, poly2 & 26 & 22 & 19 & 19 & 19 \\
64/1/8 & linear & 19 & 19 & 19 & 19 & 19 \\
128/2/16 & GELU, poly2 & 26 & 22 & 22 & 22 & 22 \\
128/2/16 & ReLU & 26 & 22 & 21 & 21 & 21 \\
128/2/16 & linear & 21 & 21 & 21 & 21 & 21 \\
\bottomrule
\end{tabular}
\end{table}

With a fill of $-100$ ($e^{-100}\approx4\times10^{-44}$, negligible in fp32), \texttt{logrows}
ranges $19$--$24$ for all variants and tracks the row count. Domain-size
comparisons between circuits are therefore only meaningful at a fixed,
modest fill; we report all grid results at $-100$. Row counts were constant
across fill values (to within $2\times10^{-5}$). This is consistent with the
observation of \citet{telescarse2025} that narrowing activation ranges shrinks
the lookup tables.

% ============================================================
\section{Quality Study}
\label{sec:quality}
% ============================================================

\paragraph{Protocol.}
WikiText-2~\citep{merity2017wikitext} (raw), GPT-2~\citep{radford2019gpt2} BPE,
non-overlapping blocks of 256 tokens (2.39M training tokens). Every model is
randomly initialized and trained with AdamW (peak LR $3\times10^{-4}$, cosine
schedule, 200 warm-up steps, weight decay $0.1$, gradient-norm clipping at $1.0$,
no dropout; the same settings for every model, with no per-variant tuning),
batch 16, 5{,}000 steps ($\approx8.6$ epochs), three seeds shared across
configurations. Validation perplexity is evaluated every
250 steps; we select the checkpoint with the lowest validation loss (the trainer's
per-batch mean, which differs from the token-weighted value reported below by
about $0.5$ perplexity) and report validation and test perplexity at that
checkpoint (token-weighted, full splits). Validation perplexity at the selected
checkpoint is therefore optimistically biased, and test perplexity is the
unbiased figure.
$\pm$ is the sample standard deviation over seeds ($n{=}3$). The reference is
pretrained GPT-2-small evaluated zero-shot on the same blocks. Except where
marked, models use rotary positions.

\begin{table}[t]
\centering
\caption{Perplexity at the validation-selected checkpoint (mean $\pm$ s.d., 3
  seeds). ``Final val'' is the last-step validation perplexity. Modelled cost is
  the (uncalibrated) analytical model at sequence length 256, relative to the
  GPT-2-small architecture. Distillation uses a frozen pretrained GPT-2 teacher
  ($\alpha{=}0.7$) and is not a from-scratch result.}
\label{tab:quality}
\resizebox{\textwidth}{!}{%
\begin{tabular}{lrrrrrr}
\toprule
Model & Params & Modelled cost & Val PPL & Test PPL & Final val & Best step \\
\midrule
GPT-2-small, pretrained, zero-shot & 124M & --- & 44.4 & 43.5 & --- & --- \\
GPT-2-small architecture, learned positions & 124M & $1.0\times$ & $179.2\pm2.2$ & $187.0\pm2.6$ & 275.3 & 2250 \\
GPT-2-small architecture, rotary & 124M & $1.0\times$ & $150.9\pm0.1$ & $158.4\pm0.9$ & 243.4 & 2250 \\
640d/10L GELU & 81M & $1.7\times$ cheaper & $147.6\pm1.7$ & $156.0\pm1.9$ & 198.5 & 2750 \\
512d/8L ReLU & 51M & $3.3\times$ cheaper & $151.7\pm0.7$ & $157.8\pm1.1$ & 204.0 & 2250 \\
384d/10L poly2 & 37M & $4.5\times$ cheaper & $155.0\pm1.2$ & $161.0\pm1.2$ & 223.4 & 2417 \\
384d/6L GELU & 30M & $7.6\times$ cheaper & $161.2\pm1.0$ & $167.0\pm1.4$ & 200.0 & 2250 \\
384d/6L ReLU & 30M & $7.6\times$ cheaper & $159.7\pm1.7$ & $165.7\pm1.7$ & 199.6 & 2250 \\
384d/6L poly2 & 30M & $7.6\times$ cheaper & $160.4\pm1.7$ & $166.6\pm1.8$ & 229.0 & 2250 \\
384d/6L linear & 30M & $7.6\times$ cheaper & $237.3\pm1.5$ & $244.0\pm1.8$ & 240.7 & 3500 \\
384d/4L linear & 26M & $11.4\times$ cheaper & $247.0\pm3.1$ & $253.6\pm2.8$ & 251.8 & 3500 \\
384d/4L linear, distilled from GPT-2 & 26M & $11.4\times$ cheaper & $185.2\pm2.0$ & $187.7\pm1.9$ & 185.3 & 4917 \\
\bottomrule
\end{tabular}}
\end{table}

\paragraph{Findings.}
(i) At matched shape (384d/6L) GELU, ReLU and poly2 are within about $1\%$ of
each other ($161.2$, $159.7$, $160.4$ validation; $167.0$, $165.7$, $166.6$ test).
Paired by seed, ReLU is the lowest in all three seeds (below GELU by $1.5\pm0.7$
and below poly2 by $0.7\pm0.5$ validation perplexity, mean $\pm$ s.d.); with
three seeds this has a one-sided sign-test $p{=}0.125$, so the ordering is
suggestive and small, and we do not claim equality or a ranking. (ii) Capacity helps: perplexity falls from $161$ (30M) to
$148$ (81M), and rotary beats learned positions by $28$ perplexity points at GPT-2
size, where the learned-position model over-fits fastest. (iii) The ``linear''
variant (no softmax, identity FFN) is $47\%$ worse than GELU at matched shape
(it may also be under-tuned, since the learning rate is shared).
(iv) Distillation lowers the 26M linear model's perplexity by $25\%$
($247.0\to185.2$; that run was still improving at step 5{,}000). (v) Nothing
trained from scratch approaches pretrained GPT-2 ($44.4$). (vi) All softmax
models over-fit: the best checkpoint is at steps $2{,}250$--$2{,}750$ and the
final-step perplexity is $18$--$61\%$ higher, so conclusions drawn from final
checkpoints would be wrong (Table~\ref{tab:quality}, ``Final val'').

\paragraph{Attention and FFN activation separately.}
Table~\ref{tab:dq} crosses the two factors at 384d/6L with the same protocol
(the identity-FFN/linear-attention cell is the ``linear'' variant above).

\begin{table}[t]
\centering
\caption{Validation perplexity at the selected checkpoint, 384d/6L (mean $\pm$
  s.d., 3 seeds). Test perplexities have the same ordering ($244.0$, $198.1$,
  $219.3$, $165.7$, $221.7$, $167.0$ in the order of the cells below, row by row,
  linear then softmax).}
\label{tab:dq}
\begin{tabular}{lcc}
\toprule
FFN activation & Linear attention & Softmax attention \\
\midrule
identity & $237.3\pm1.5$ & $192.4\pm0.6$ \\
ReLU     & $211.8\pm2.5$ & $159.7\pm1.7$ \\
GELU     & $214.7\pm1.5$ & $161.2\pm1.0$ \\
\bottomrule
\end{tabular}
\end{table}

Both factors matter. Replacing softmax by linear attention raises perplexity by
$23\%$ with an identity FFN and by $33\%$ with ReLU or GELU; removing the FFN
nonlinearity raises it by $19$--$20\%$ with softmax attention and by $11$--$12\%$
with linear attention. The $47\%$ gap of the ``linear'' variant to GELU is
therefore a combination of both effects, not attributable to either alone.

% ============================================================
\section{Reading Cost and Quality Together}
\label{sec:joint}
% ============================================================

Table~\ref{tab:joint} places the two measurements side by side for the five
variants for which both exist. Rows are the measured EZKL rows from
Table~\ref{tab:decomp} (width 128, $\nlayers{=}2$, $s{=}16$), relative to the
linear-attention/identity-FFN variant; perplexity is from the 384d/6L study. The
two were measured on different shapes, so the table supports only the
\emph{ordering} of variants, not absolute costs at the quality-study scale.

\begin{table}[t]
\centering
\caption{Measured relative rows (width 128, 2 layers, sequence 16) against
  validation perplexity (384d/6L) for five variants.}
\label{tab:joint}
\begin{tabular}{lcc}
\toprule
Variant (attention / FFN) & Relative rows & Val PPL \\
\midrule
linear / identity & 1.00 & $237.3\pm1.5$ \\
softmax / identity & 1.02 & $192.4\pm0.6$ \\
softmax / ReLU & 1.09 & $159.7\pm1.7$ \\
softmax / poly2 & 1.23 & $160.4\pm1.7$ \\
softmax / GELU & 1.44 & $161.2\pm1.0$ \\
\bottomrule
\end{tabular}
\end{table}

Three variants are non-dominated on rows and perplexity: linear/identity,
softmax/identity and softmax/ReLU. Softmax attention costs $2\%$ more rows than
linear attention at this width ($4$--$24\%$ at width 64 as $s$ grows from 8 to 64,
Table~\ref{tab:seqT}) and lowers perplexity by $19\%$ (identity FFN); a ReLU FFN
costs $7\%$ more and lowers it by another $17\%$ ($192.4\to159.7$). GELU and poly2
use $1.44\times$ and $1.23\times$ the base rows ($32\%$ and $12\%$ more than ReLU)
with no perplexity gain. Proving time does not carry that separation
(Section~\ref{sec:proving}): nonlinear variants take $20$--$30\%$ longer to prove
than the linear variant, but ReLU, poly2 and GELU trade places across shapes
within the run-to-run noise. We therefore conclude only that ReLU has the fewest
rows among the three nonlinearities and, in this protocol, equal perplexity; we
cannot show that this yields a proving-time advantage. These are statements about
this grid and this training protocol only (Section~\ref{sec:limits}).

% ============================================================
\section{Implementation Notes and Sensitivity}
\label{sec:causal}
% ============================================================

\paragraph{Causal masking.}
Attention is causal: a lower-triangular mask is applied to the scores, and a
regression test asserts, for every activation, with rotary on and off and for
every attention/FFN combination used here, that changing a future token leaves
earlier-position logits unchanged to $10^{-6}$. The mask matters for correctness:
without it, next-token targets leak and perplexities are not comparable with those
of causal language models (we found and fixed this in our own code before running
any experiment reported here).

\paragraph{Checkpoint selection.}
Validation-selected checkpoints and final-step perplexities differ by $18$--$61\%$
(Table~\ref{tab:quality}), enough to change two conclusions: final-step
perplexity would make poly2 appear $14\%$ worse than GELU and make capacity appear
not to matter. We report the selected checkpoint and keep final-step numbers.

\paragraph{Mask fill value.}
Section~\ref{sec:domain} gives the sensitivity of \texttt{logrows} to the mask
fill; domain sizes are comparable only at a fixed, modest fill.

% ============================================================
\section{Discussion and Limitations}
\label{sec:limits}
% ============================================================

\paragraph{Takeaways.}
For small causal transformers compiled and proved with EZKL~23.0.5: (1) an
analytical constraint model is not better than the product of the dimensions for
ranking, is uncalibrated, and extrapolates worse than fitted alternatives (at width 256--512 its fitted
prediction errs by $122\%$ against $10\%$ for the dimension product); a
calibrated additive model is best on held-out grid values but not out of range
(it puts no weight on the matrix-product term, an exploratory result given
collinear features); (2) at short sequences the FFN nonlinearity drives the row
overhead for GELU and poly2 (ReLU's is comparable to softmax's at width 64), but softmax grows with sequence length and at
width 64 matches GELU's overhead by $s{=}64$; (3) compiled rows predict proving time across shapes (Spearman $0.89$)
through a power-of-two domain that can double the cost, but repeated runs differ by up
to $1.9\times$ on a shared machine and do not separate the nonlinearities; (4) ReLU,
poly2 and GELU are within $1\%$ in perplexity (ReLU nominally lowest), and ReLU has
the fewest rows among them;
(5) domain-size comparisons are only meaningful at a fixed,
modest mask fill.

\paragraph{Limitations.}
\begingroup\sloppy
(1) \emph{Scale:} the compiled grid has width $\le128$, depth $\le4$, sequence
length $\le32$ (decomposition to $64$) and vocabulary 512, and the out-of-range
check reaches width 512 with 1--2 layers and sequence length $\le64$; proofs are limited to
width $\le128$, 1--2 layers, sequence length $\le16$ and domain $\le2^{21}$. We do
not measure GPT-2-scale circuits, the regime where the analytical argument is
strongest, and the fitted laws should not be extrapolated beyond the range tested.
(2) \emph{Proving measurements are noisy:} two runs for the 20 configurations with domain at most
$2^{20}$ and one for the four at $2^{21}$, on a machine shared with other jobs
(load average $20$--$40$; run-to-run time ratios $0.82$--$1.91\times$), with a locally
generated, insecure reference string; we could not obtain an idle machine.
(3) \emph{Single prover and version:}
EZKL~23.0.5, random weights and fixed random calibration tokens; the fixed-point
scale was chosen by calibration (replicated with the scale pinned); other
backends are not measured. (4) \emph{Quality and cost are measured on different
models:} the quality-study models were not compiled or proved, and the
modelled-cost column of Table~\ref{tab:quality} comes from the uncalibrated
model; Table~\ref{tab:joint} joins different shapes and supports ordering only.
(5) \emph{Quality data:} WikiText-2 is small, models over-fit within about 2{,}500
steps, a single learning rate and no dropout or regularization sweep are used, and
$n{=}3$ seeds, so differences of about $1\%$ (such as among GELU, ReLU and poly2)
are consistent in sign but small, and activation rankings could change under a
tuned learning rate or regularization; the ``linear'' variants may be under-tuned.
(6) \emph{Quantized accuracy} of the compiled circuits is not measured.
(7) \emph{Implementation confound:} linear attention multiplies the scores by the
causal mask while softmax attention uses a masked fill followed by softmax, so
the measured softmax overhead also includes that implementation difference. (8) \emph{No search
procedure} is evaluated; the repository contains search code about which this
paper makes no claim.
\endgroup

\paragraph{Reproducibility.}
\begingroup\sloppy
Code, the 108 EZKL records at calibrated and pinned scale, the decomposition
records (sequence lengths 8--64), the 16 out-of-range records, the mask-fill
sweep, the proof records (both runs) and
all training-run records (with learning curves) are in
\url{https://github.com/npow/ezkl-transformer-cost-study} (release tag
\path{v1}; scripts \path{run_ezkl_validation.py},
\path{run_ezkl_proving.py} (repeat runs via its \path{--out} option),
\path{run_maskfill_sweep.py}, \path{run_causal_lm_study.py},
\path{analyze_validation.py}; raw data under \path{artifacts/}; the cost model's
operator table is vendored in \path{proxy/op_costs.py}). All reported numbers are
produced by those scripts.
\endgroup

\paragraph{Conclusion.}
Within the range we could measure and prove, an analytical constraint model adds
little over the product of the dimensions and fitted alternatives beat it in
held-out and out-of-range row prediction; the FFN nonlinearity is the main
activation-dependent row cost at short sequences while softmax grows with
sequence length; compiled rows predict proving time through a power-of-two
domain that can double the cost, although proving times on a shared machine are too
noisy to separate ReLU, poly2 and GELU; and ReLU has the fewest rows among the
nonlinearities at no perplexity penalty in this protocol. Whether any of this holds
at GPT-2 scale remains to be measured with full-size circuits.

\bibliographystyle{plainnat}

\begin{thebibliography}{99}

\bibitem[Bian et~al.(2020)]{nass2020}
Bian, S., Jiang, W., Lu, Q., Shi, Y., and Sato, T.
\newblock {NASS}: Optimizing secure inference via neural architecture search.
\newblock In \emph{European Conference on Artificial Intelligence (ECAI)}, 2020.
  arXiv:2001.11854.

\bibitem[Cai et~al.(2019)]{cai2019proxylessnas}
Cai, H., Zhu, L., and Han, S.
\newblock {ProxylessNAS}: Direct neural architecture search on target task and
  hardware.
\newblock In \emph{International Conference on Learning Representations
  (ICLR)}, 2019.

\bibitem[Zkonduit(2024)]{ezkl}
Zkonduit.
\newblock {EZKL}: a library and command-line tool for zero-knowledge machine
  learning inference (version 23.0.5).
\newblock \url{https://github.com/zkonduit/ezkl}.

\bibitem[Husistein et~al.(2025)]{husistein2025near}
Husistein, R.~T., Reiher, M., and Eckhoff, M.
\newblock {NEAR}: A training-free pre-estimator of machine learning model
  performance.
\newblock In \emph{International Conference on Learning Representations
  (ICLR)}, 2025. arXiv:2408.08776.

\bibitem[Javaheripi et~al.(2022)]{javaheripi2022lts}
Javaheripi, M., de~Rosa, G.~H., Mukherjee, S., Shah, S., Religa, T.~L.,
  Mendes, C.~C.~T., Bubeck, S., Koushanfar, F., and Dey, D.
\newblock {LiteTransformerSearch}: Training-free neural architecture search for
  efficient language models.
\newblock In \emph{Advances in Neural Information Processing Systems
  (NeurIPS)}, 2022. arXiv:2203.02094.

\bibitem[Kang et~al.(2022)]{kang2022scaling}
Kang, D., Hashimoto, T., Stoica, I., and Sun, Y.
\newblock Scaling up trustless {DNN} inference with zero-knowledge proofs.
\newblock arXiv preprint arXiv:2210.08674, 2022.

\bibitem[Lee et~al.(2019)]{lee2019snip}
Lee, N., Ajanthan, T., and Torr, P.~H.~S.
\newblock {SNIP}: Single-shot network pruning based on connection sensitivity.
\newblock In \emph{International Conference on Learning Representations
  (ICLR)}, 2019.

\bibitem[Maheri et~al.(2025)]{telescarse2025}
Maheri, M.~M., Haddadi, H., and Davidson, A.
\newblock {TeleSparse}: Practical privacy-preserving verification of deep neural
  networks.
\newblock \emph{Proceedings on Privacy Enhancing Technologies (PoPETs)},
  2025(4):861--880, 2025.

\bibitem[Merity et~al.(2017)]{merity2017wikitext}
Merity, S., Xiong, C., Bradbury, J., and Socher, R.
\newblock Pointer sentinel mixture models.
\newblock In \emph{International Conference on Learning Representations
  (ICLR)}, 2017.

\bibitem[Qu et~al.(2025)]{zkgpt2025}
Qu, W., Sun, Y., et~al.
\newblock {zkGPT}: An efficient non-interactive zero-knowledge proof framework
  for {LLM} inference.
\newblock In \emph{USENIX Security Symposium}, pages 2045--2063, 2025.

\bibitem[Radford et~al.(2019)]{radford2019gpt2}
Radford, A., Wu, J., Child, R., Luan, D., Amodei, D., and Sutskever, I.
\newblock Language models are unsupervised multitask learners.
\newblock OpenAI Blog, 2019.

\bibitem[Serianni and Kalita(2023)]{serianni2023understanding}
Serianni, A.\ and Kalita, J.
\newblock Training-free neural architecture search for {RNNs} and transformers.
\newblock In \emph{Annual Meeting of the Association for Computational
  Linguistics (ACL)}, pages 2522--2540, 2023.

\bibitem[Sun et~al.(2024)]{zkllm2024}
Sun, H., Li, J., and Zhang, H.
\newblock {zkLLM}: Zero knowledge proofs for large language models.
\newblock In \emph{ACM Conference on Computer and Communications Security
  (CCS)}, 2024. arXiv:2404.16109.

\bibitem[Tan et~al.(2019)]{tan2019mnasnet}
Tan, M., Chen, B., Pang, R., Vasudevan, V., Sandler, M., Howard, A., and
  Le, Q.~V.
\newblock {MnasNet}: Platform-aware neural architecture search for mobile.
\newblock In \emph{IEEE/CVF Conference on Computer Vision and Pattern
  Recognition (CVPR)}, 2019.

\bibitem[Wang(2026)]{nanozk2026}
Wang, Z.~G.
\newblock {NANOZK}: Layerwise zero-knowledge proofs for verifiable large
  language model inference.
\newblock arXiv preprint arXiv:2603.18046, 2026.

\bibitem[Wu et~al.(2019)]{wu2019fbnet}
Wu, B., Dai, X., Zhang, P., Wang, Y., Sun, F., Wu, Y., Tian, Y., Vajda, P.,
  Jia, Y., and Keutzer, K.
\newblock {FBNet}: Hardware-aware efficient {ConvNet} design via differentiable
  neural architecture search.
\newblock In \emph{IEEE/CVF Conference on Computer Vision and Pattern
  Recognition (CVPR)}, 2019.

\end{thebibliography}

\end{document}
